% ==============================================================================
% NOTULA™ CLI OFFICIAL REFERENCE MANUAL & TECHNICAL SPECIFICATION
% Ideazione & Sviluppo: Ing. Mario Fantini
% Sito Ufficiale: https://mariofantini.eu
% Versione Specifica: 2.2.0 (Agosto 2026)
% ==============================================================================
\documentclass[11pt,a4paper]{article}

\usepackage[utf8]{inputenc}
\usepackage[T1]{fontenc}
\usepackage[italian]{babel}
\usepackage{geometry}
\geometry{a4paper, top=2.5cm, bottom=2.5cm, left=2.5cm, right=2.5cm}
\usepackage{hyperref}
\hypersetup{
    colorlinks=true,
    linkcolor=blue!80!black,
    urlcolor=blue!80!black,
    citecolor=blue!80!black,
    pdftitle={Notula CLI Reference Manual - Ing. Mario Fantini},
    pdfauthor={Ing. Mario Fantini}
}
\usepackage{listings}
\usepackage{xcolor}
\usepackage{tcolorbox}
\usepackage{tabularx}
\usepackage{booktabs}
\usepackage{amsmath,amssymb}
\usepackage{fancyhdr}

\pagestyle{fancy}
\fancyhf{}
\fancyhead[L]{\textbf{Notula™ CLI Manual} -- v2.2.0}
\fancyhead[R]{Ing. Mario Fantini}
\fancyfoot[C]{\thepage}

\definecolor{cliDark}{RGB}{13, 17, 23}
\definecolor{cliCyan}{RGB}{56, 189, 248}
\definecolor{cliGreen}{RGB}{52, 211, 153}
\definecolor{cliYellow}{RGB}{250, 204, 21}

\lstdefinestyle{notulacli}{
    backgroundcolor=\color{cliDark},
    basicstyle=\ttfamily\small\color{white},
    keywordstyle=\color{cliGreen}\bfseries,
    stringstyle=\color{cliYellow},
    commentstyle=\color{gray},
    breaklines=true,
    frame=single,
    rulecolor=\color{gray!50},
    numbers=none
}

\begin{document}

\begin{titlepage}
    \centering
    \vspace*{1.5cm}
    {\huge\bfseries NOTULA™ ENTERPRISE CLI}\\[0.4cm]
    {\Large Manuale Ufficiale dei Comandi \& Specifica Tecnica di Sistema}\\[0.2cm]
    {\large Versione 2.2.0 -- Edizione Ufficiale Agosto 2026}\\[2cm]
    
    \begin{tcolorbox}[colback=blue!5!white,colframe=blue!75!black,title=\textbf{Informazioni sull'Opera \& Proprietà Intellettuale}]
        \textbf{Autore \& Progettista:} Ing. Mario Fantini\\
        \textbf{Sito Web Ufficiale:} \href{https://mariofantini.eu}{https://mariofantini.eu}\\
        \textbf{Architettura Software:} Notula Memo Engine Enterprise v2.2\\
        \textbf{Licenza \& Diritti:} Proprietary \& Confidential -- Tutti i Diritti Riservati.
    \end{tcolorbox}
    
    \vfill
    {\small Documento ad uso tecnico, ingegneristico e professionale.\\
    Generato e compilato con standard accademico \LaTeX.}
\end{titlepage}

\tableofcontents
\newpage

\section{Introduzione \& Architettura dell'Interprete CLI}
L'interprete \textbf{Notula™ CLI} fornisce una shell professionale basata su riga di comando accessibile sia tramite la scorciatoia globale \texttt{Ctrl+Shift+P} sia dal pulsante Terminale presente nella barra superiore dell'applicazione.

\subsection{Caratteristiche Ingegneristiche}
\begin{itemize}
    \item \textbf{Sincronizzazione di Gruppo (\texttt{groupID}):} Ogni promemoria generato all'interno di una serie ricorrente eredita il medesimo identificativo di gruppo univoco. Modificando titolo, note o crittografia tramite il comando \texttt{edit}, le variazioni vengono propagate a cascata preservando le date individuali di calendario.
    \item \textbf{Crittografia a 3 Strati:} Supporto nativo hardware \texttt{AES-GCM-256} derivata da Master Passphrase con 100.000 iterazioni \texttt{PBKDF2}.
    \item \textbf{Interoperabilità Totale:} Esportazione ed importazione nei 6 formati universali: JSON, XML, Markdown, iCalendar (ICS), PDF A4 Ink-Friendly e Plain Text.
\end{itemize}

\section{Sintassi Generale \& Notazione}
Tutti i comandi rispettano la notazione standard POSIX/GNU:
\begin{lstlisting}[style=notulacli]
notula> <comando> [--opzione <valore>] [--flag]
\end{lstlisting}

\section{Riferimento Completo dei Comandi}

\subsection{1. Comando \texttt{add} -- Creazione Promemoria}
\textbf{Sintassi:} \texttt{add --title <str> --date <YYYY-MM-DD> [opzioni...]}
\begin{lstlisting}[style=notulacli]
notula> add --title "Udienza Civile" --date 2026-09-20
notula> add --title "Hosting VPS" --date 2026-09-01 --repeat monthly --encrypt
\end{lstlisting}

\subsection{2. Comando \texttt{edit} -- Modifica \& Cascading Sync}
\textbf{Sintassi:} \texttt{edit --id <id> [--title <str>] [--date <YYYY-MM-DD>] [--desc <str>] [--security <0|1|2>] [--encrypt]}
\begin{lstlisting}[style=notulacli]
notula> edit --id n_17240012345 --title "Udienza Rinviata" --date 2026-10-15
\end{lstlisting}

\subsection{3. Comando \texttt{export} -- Esportazione Multiformato}
\textbf{Sintassi:} \texttt{export [--id <id> | --all | --year <YYYY>] --format <json|xml|md|ics|pdf|txt>}
\begin{lstlisting}[style=notulacli]
notula> export --all --format json
notula> export --year 2026 --format ics
notula> export --id n_17240012345 --format pdf
\end{lstlisting}

\subsection{4. Comando \texttt{pdf} -- Generazione Documentale Ink-Friendly}
\textbf{Sintassi:} \texttt{pdf [--id <id>]}
\begin{lstlisting}[style=notulacli]
notula> pdf --id n_17240012345
\end{lstlisting}

\subsection{5. Comando \texttt{import} -- Iniezione Dati}
\textbf{Sintassi:} \texttt{import [--json <raw_string> | --xml <raw_string>]}
\begin{lstlisting}[style=notulacli]
notula> import --json '[{"title":"Scadenza","expirationDate":"2026-11-30"}]'
notula> import
\end{lstlisting}

\subsection{6. Comando \texttt{rm} -- 9 Modalità di Eliminazione Mirata}
\textbf{Sintassi:} \texttt{rm [--id <id>] [--group <groupID>] [--title <query>] [--expired] [--all] [--force]}
\begin{lstlisting}[style=notulacli]
notula> rm --id n_17240012345
notula> rm --group grp_17240012345_abc
notula> rm --expired
notula> rm --all --force
\end{lstlisting}

\subsection{7. Comando \texttt{passwd} -- Sicurezza \& Master Passphrase}
\textbf{Sintassi:} \texttt{passwd <nuova_passphrase>}
\begin{lstlisting}[style=notulacli]
notula> passwd MiaNuovaPasswordSicura2026!
\end{lstlisting}

\subsection{8. Comandi di Query \& Diagnostica: \texttt{ls}, \texttt{show}, \texttt{sync}, \texttt{stats}}
\begin{lstlisting}[style=notulacli]
notula> ls --year 2026 --expired
notula> show n_17240012345
notula> sync
notula> stats
\end{lstlisting}

\section{Licenza \& Contatti}
\textbf{Notula™ Memo Engine} è un'opera di ingegno ideata e sviluppata dall'\textbf{Ing. Mario Fantini}.\\
Per documentazione, licenze personalizzate e supporto:\\
\href{https://mariofantini.eu}{https://mariofantini.eu} -- \texttt{marfant7@gmail.com}

\end{document}
